\section{Output Validation Using VAL}

After the benchmark has obtained a textual answer from an LLM and extracted the
candidate PDDL actions, the plan still has to be checked against the formal
domain and problem definition. This is the role of VAL in the framework. The
parser can say whether the model produced something that looks like an action
sequence; VAL decides whether that sequence is actually executable and whether
it reaches the goal.

\subsection{Role of VAL in the Benchmark}

VAL is used as an external symbolic validator. For each benchmark case, the
framework provides VAL with three inputs:

\begin{itemize}
    \item the PDDL domain file, containing predicates, actions, preconditions,
    and effects;
    \item the PDDL problem file, containing objects, initial state, and goal;
    \item the candidate plan extracted from the LLM output.
\end{itemize}

This separation is important because the benchmark does not rely on textual
judgement to decide whether a plan is correct. A model can produce a plausible
or well explained answer, but the answer is counted as solved only if the final
action sequence validates against the corresponding PDDL task.

\subsection{Validator Setup and Preflight}

VAL is not a Python dependency of the project. It is an external executable that
can be resolved in three ways: by putting \texttt{Validate} or
\texttt{Validate.exe} on the system path, by passing an explicit executable path
with \texttt{--validator-command}, or by using one of the bundled platform
utilities under \texttt{utils/linux64/bin/} or \texttt{utils/win64/bin/}.

The benchmark can also run a task preflight before launching model generation.
In this stage the framework checks the selected domain and problem files with
VAL. This prevents malformed or unsupported PDDL inputs from being interpreted
later as model failures. If VAL cannot be resolved, or if the preflight fails,
the error is stored as an orchestration problem instead of being mixed with the
LLM planning results.

\subsection{Validation Flow}

The validation path is the same for every model backend. The backend changes how
the answer is generated, but not how the answer is checked:

\begin{center}
\texttt{LLM output}
$\rightarrow$ \texttt{parsed actions}
$\rightarrow$ \texttt{temporary plan file}
$\rightarrow$ \texttt{VAL(domain, problem, plan)}
$\rightarrow$ \texttt{ValidationResult}
$\rightarrow$ \texttt{scored artifact}
\end{center}

The temporary plan file is created only to invoke VAL with the expected command
line interface. After validation, the framework normalizes the validator output
into a common result schema. This makes the rest of the benchmark independent
from the exact textual format printed by the external validator.

\subsection{Normalized Validation Result}

The normalized validation record stores both the final decision and the
diagnostic information needed for analysis and repair.

\begin{table}[H]
\centering
\renewcommand{\arraystretch}{1.2}
\begin{tabularx}{\textwidth}{lX}
\toprule
\textbf{Field} & \textbf{Meaning} \\
\midrule
\texttt{valid} & Boolean result used to decide whether the candidate plan solved the task. \\
\texttt{status} & Normalized status such as \texttt{valid}, \texttt{invalid}, \texttt{parse\_error}, \texttt{timeout}, or \texttt{validator\_error}. \\
\texttt{error\_type} & Compact failure class used by metrics and repair feedback. \\
\texttt{failed\_step} & First failing step reported by VAL, when available. \\
\texttt{failed\_action} & Problematic action extracted from the validator output, when available. \\
\texttt{goal\_satisfied} & Whether the final state satisfies the goal. \\
\texttt{plan\_length} & Number of action lines submitted to validation. \\
\texttt{validation\_time\_ms} & Time spent in the external validator call. \\
\texttt{raw\_validator\_output} & Original VAL output kept for debugging and later inspection. \\
\bottomrule
\end{tabularx}
\caption{Main fields stored in the normalized validation result.}
\label{tab:validation-result-fields}
\end{table}

\subsection{Prefix Validation}

The framework does not validate only the full action sequence. It validates
every non-empty prefix of the parsed final answer. This provides a more precise
view of the model behaviour.

For example, a model may reach the goal after five actions and then continue
with unnecessary or invalid actions. In that case, validating only the complete
sequence would mark the attempt as invalid, while prefix validation reveals that
a valid solution was present inside the answer. The scored artifact therefore
records the first valid prefix length, the first valid plan text, whether the
final full plan is valid, and how many extra actions appear after the first
valid prefix.

This distinction is useful for the later reasoning analysis: it separates models
that never find a valid plan from models that find a plan but fail to stop cleanly
after satisfying the goal.

\subsection{Failure Types and Repair Feedback}

VAL failures are mapped into a small shared taxonomy. The most relevant classes
are:

\begin{itemize}
    \item \texttt{invalid\_precondition}: an action is not applicable in the
    current state;
    \item \texttt{unsatisfied\_goal}: the sequence executes but does not reach
    the required goal;
    \item \texttt{unknown\_action}: the plan contains an action not defined in
    the domain;
    \item \texttt{syntax\_error}: the plan format is not accepted by the
    validator;
    \item \texttt{timeout}, \texttt{validator\_unavailable}, and
    \texttt{validator\_crash}: technical validation failures, kept separate
    from planning mistakes.
\end{itemize}

In the iterative repair protocol, this normalized failure information is turned
into feedback for the next model attempt. The feedback can include the validation
status, the error type, the failed step, the failed action, and a compact
description of the validator output. This makes the repair loop grounded in
symbolic validation rather than in a generic request to try again.

\subsection{Reasoning Validation}

Some models expose a separate reasoning trace in addition to the final answer.
The benchmark keeps the final answer as the official plan used for scoring, but
it can also validate an action sequence extracted from the reasoning text. This
second validation is diagnostic: it helps identify cases where the model derives
a valid plan in its reasoning but fails to transfer it correctly into the final
answer.

This distinction avoids giving credit for hidden reasoning when the final answer
is invalid, while still preserving useful information for the reasoning
capability analysis in the following section.

\subsection{Metrics Derived from Validation}

The validation result is the source for the core execution metrics. A case is
considered solved only when the normalized validation result is valid. From this
record the framework derives \texttt{validity\_at\_1},
\texttt{validity\_at\_k}, \texttt{repair\_success},
\texttt{iterations\_to\_valid}, \texttt{plan\_length},
\texttt{error\_type}, and whether the run hit the iteration limit.

Therefore, VAL is not only a final correctness check. It is the component that
connects raw model generations to comparable benchmark scores.

\subsection{Limitations}

The validation layer deliberately checks formal plan correctness, not the
quality of the natural-language reasoning. It also depends on the availability
and behaviour of an external executable, so timeout and validator errors are
reported separately from genuine planning failures. Finally, the parser that
normalizes VAL output uses conservative textual heuristics; for this reason the
raw validator output is stored together with the normalized fields whenever
manual inspection is needed.
